%%%
%%% OpenCAEシンポジウムTeXテンプレートファイル
%%% template_OpenCAE_symposium.tex
%%% OpenCAEシンポジウム2018版
%%%
%%
%% ltjocはOpenCAE論文集・シンポジウム用のクラスファイルです．変更しないでください．
%% 本文が英語の場合には，オプションにenglishを指定してください．
\documentclass{../../style/ltjoc}
%\documentclass[english]{../../style/ltjoc}
%%
%% 表題(title)，副題(subtitle)，著者(author)，所属(affiliation)
%% 本文が英語の場合には，こちらに英文表題，英文副題，英文著者，英文所属を記述します.
\title{オープンCAEシンポジウム2018印刷用原稿の書き方（和文表題）}
% 副題が無い場合にはコメントアウトします．
\subtitle{\TeX のテンプレート（和文副題）} 
\author{%
氏 名$^{1\dagger}$%
\hspace{1\zw}%
氏 名$^{2}$%
\hspace{1\zw}%
氏 名$^{3}$%
}
\affiliation{%
${}^{1}$所属%
\hspace{1\zw}%
${}^{2}$所属%
\hspace{1\zw}%
${}^{3}$所属%
} 
%%
%% Corresponding authorの電子メールアドレス
%% 本論文について連絡が取れる著者の電子メールアドレスを記載してください．
\AuthorsEmail{office@opencae.or.jp}
%%
%% 英文表題(etitle)，英文副題(esubtitle)，英文著者(eauthor)，英文所属(eaffiliation)
%% 本文が英語の場合には表示されません．
\etitle{Manuscripts preparation guide for the Open CAE symposium 2018 (English Title)}
% 副題が無い場合にはコメントアウトします．
\esubtitle{The case of \TeX (English Sub-Title)}
\eauthor{%
Firstname FAMILYNAME$^{*\dagger}$%
\hspace{1em}%
Firstname FAMILYNAME$^{**}$%
\hspace{1em}%
Firstname FAMILYNAME$^{***}$%
}
\eaffiliation{%
${}^{*}$affiliation%
\hspace{1em}%
${}^{**}$affiliation%
\hspace{1em}%
${}^{***}$affiliation%
}
%%
%% キーワード
\keywords{Keyword1, Keyword2, Keyword3, Keyword4, Keyword5}
%%
%% 英文概要
%% 英文概要を省略する場合には，abstract環境の定義をしないでください．
\begin{abstract}
The quick brown fox jumps over the lazy dog.
The quick brown fox jumps over the lazy dog.
The quick brown fox jumps over the lazy dog.
The quick brown fox jumps over the lazy dog.
The quick brown fox jumps over the lazy dog.
The quick brown fox jumps over the lazy dog.
\end{abstract}
%%
%% luatexja-fontspecパッケージ
\usepackage{luatexja-fontspec}
\defaultfontfeatures{Ligatures=TeX}
%% luatexja-presetパッケージ
%% 和文フォントのプリセット設定
%% オプションにnoembed(非埋込)を指定しないでください．
%\usepackage[ipaex]{luatexja-preset} % IPAex(デフォルト)
%\usepackage[ms]{luatexja-preset} % MS
%\usepackage[hiragino-pro]{luatexja-preset} % ヒラギノPro
%%
%% 欧文フォントの指定
%% 使用できるフォントについては，以下のコマンドで調べてください．
%% $ luaotfload-tool --list=*
%%
%% 欧文通常フォント
%\setmainfont{Cambria} % Cambria
%\setmainfont{Times New Roman} % Times New Roman
%\setmainfont{TeXGyreTermes} % TeXGyreTermes
%%
%% 欧文Sans-serifフォント
%\setsansfont{Calibri} % Calibri
%\setsansfont{Arial} % Arial
%\setsansfont{Helvetica} % Helvetica
%\setsansfont{TeXGyreHeros} % TeXGyreHeros
%%
%% 欧文monospaceフォント
%\setmonofont{Consolas} % Consolas
%\setmonofont{Courier New} % Courier New
%\setmonofont{Lucida Console} % Lucida Console
%%
%% subfigureパッケージ
\usepackage{subfigure}
%%
%% graphicxパッケージ
\usepackage{graphicx}
%%
%% hyperrefパッケージ
\usepackage[
 pdfencoding=auto,
 bookmarks=true,
 bookmarksnumbered=true,
 colorlinks=true,
 allcolors={blue}
]{hyperref}
%%
%% listingsパッケージ
\usepackage{listings}
\renewcommand{\lstlistingname}{Code}
\lstset{
basicstyle={\footnotesize\ttfamily},
commentstyle=\color{blue},
frame={tb},
breaklines=true,
columns=[l]{fullflexible},
numbers=left,
numberstyle={\footnotesize},
keepspaces=true
}
%%
%% autorefでの図表の参照名の再定義
\makeatletter
\if@english
  \renewcommand*{\figureautorefname}{\figurename}
  \renewcommand*{\tableautorefname}{\tablename}
\else
  \renewcommand*{\figureautorefname}{図}
  \renewcommand*{\tableautorefname}{表}
\fi
\makeatother
%%
%% ヘッダ右の設定
%% 変更しないでください．
\markright{Open CAE Symposium 2018, Dec. 7-8, 2018, Kawasaki} % Do not edit this line
%%
%% 本文
\begin{document}
%%
%% 題目などの出力
\maketitle
%%%
\section{原稿について}
%%
\subsection{様式}
\begin{itemize}
\item 用紙 : A4です．
\item 余白 : 上下25mm，左右20mmです．
  英文概要とキーワードの行は，さらに左右6.5mmの余白を設けてください．
\item 段組 : 1段組です．
\item 頁数 : 最小1ページ，最大10ページ，標準2〜4ページとします．
\end{itemize}
%%
\subsection{本文}
\begin{itemize}
\item 言語 : 日本語または英語です．
\item 句読点 : 本文が日本語の場合，句読点として，
  全角の読点「，」（カンマ）と句点「．」（ピリオド）を用いてください． 
\item 段落最初の行の字下げ : 1文字です．
\item 題目・所属
  \begin{itemize}
  \item 
    本文が日本語の場合，日本語での題目，著者名と所属の記載に続けて，
    英語でも同内容を記載してください．
  \item 
    本文が英語の場合，英語の題目および著者名と所属のみ記載ください．
  \end{itemize}
\item 英文概要 : 40から80ワード程度で記載ください．
  なお，英文概要は省略しても構いません．
\item キーワード : 講演内容を良く表すキーワードを最低3語，通常5語程度選定し，英語で記入ください.
\item 字体
  \begin{itemize}
  \item 本文 :
    明朝体・Serif系（Cambria，Times New Roman，Latin Modernなど）を使用してください．
  \item 題目・著者名・所属・見出し・図表キャプション : 
    ゴシック体・Sans-serif系（Calibri，Arial，Latin Modern Sansなど）を使用してください．
    キーワードの見出し（Keywords）はSans-serifのイタリック体を使用してください．
  \end{itemize}
\item 文字の大きさ
  \begin{itemize}
  \item 題目 : 14ポイント
  \item 副題・大見出し : 12ポイント
  \item ヘッダ・フッタ・付録におけるコード引用 : 8ポイント
  \item 上記以外（本文・著者名・所属・Abstract・キーワード・
      図表キャプション・大見出し以外の見出し・ページ番号など） : 10ポイント
  \end{itemize}
\item 参考文献
  \begin{itemize}
  \item 
  参考文献は，本文中の引用箇所の末尾に角括弧をつけた番号で表し，
  本文の末尾にまとめて列記してください\cite{Nishi:OpenCAE2018-001,BB21459576}．
  また，URLにはアクセス日付を明記してください\cite{SIST02-200}．
  \end{itemize}
\end{itemize}
%%
\subsection{図表}
\begin{itemize}
\item 言語 : 図表中の記号およびキャプションは英語とします．
\item 字体・文字の大きさ : 
  明瞭である限り，図表中の字体や文字の大きさは任意です．
\item 本文中での参照 : \autoref{fig:Logo}，
  \autoref{tab:Secretariat}などと記載してください．
\end{itemize}
%
\begin{figure}[htbp]
\centering
\subfigure[Color]{
\includegraphics[width=0.1\textwidth]{fig/Logo-c3-1.jpg}
\label{fig:Logo-c3-1}
}
\hspace{0.1\textwidth}
\subfigure[Mono.]{
\includegraphics[width=0.1\textwidth]{fig/Logo-b3.jpg}
\label{fig:Logo-b3}
}
\caption{Logo of the Open CAE Society Japan.}
\label{fig:Logo}
\end{figure}
%%%
\section{原稿提出について}
%%
\begin{itemize}
\item ファイル形式 : PDF形式で提出ください．
\item ファイルサイズ : 最大20MBです．
\item フォント : 提出されるPDFファイルに全てのフォントが埋め込まれている事を確認ください．
\item URL : リンクを付加してください．
  また，可能であれば，図表の参照についてもリンクを付加してください．
  リンクされた文字は青色としてください．
\item 提出方法 : オープンCAEシンポジウム2018のWEBページ
  \href
  {http://www.opencae.or.jp/activity/symposium/opencae_symposium2018/}
  {http://www.opencae.or.jp/activity/symposium/opencae\_symposium2018/}
  をご参照の上，提出してください．
\end{itemize}
%%%
\section{お問い合わせ}
ご不明点については，
お手数ですが，
\autoref{tab:Secretariat}のシンポジウム事務局まで，
電子メールでお問い合わせください．
\begin{table}[htbp]
  \centering
  \begin{tabular}{l|l}
    \hline
    E-mail address & symposium2018@opencae.or.jp\\
    \hline
    \end{tabular}
  \caption{Secretariat of the Open CAE symposium 2018}
  \label{tab:Secretariat}
\end{table}
%%
%% 参考文献
%% thebibliography環境を用いる場合
%% template_OpenCAE_symposium.bblに直接参考文献を記述し，
%% 以下の行を有効にしてください．
\documentclass{ltjoc}
%
\usepackage{subfigure}
\usepackage{graphicx}
\usepackage[
 pdfencoding=auto,
 bookmarks=true,
 bookmarksnumbered=true,
 colorlinks=true,
 urlcolor={black},
 linkcolor={black},
 filecolor={black},
 citecolor={black},
 allcolors={black}
]{hyperref}
\renewcommand*{\figureautorefname}{図}
\renewcommand*{\tableautorefname}{表}
% フォント
\usepackage{luatexja-otf}
\usepackage[ms]{luatexja-preset} % MSフォント
%\usepackage[ipa]{luatexja-preset} % IPAフォント
%\usepackage[hiragino-pro]{luatexja-preset} % ヒラギノフォント
%
\markrightheadings{OpenCAE Symposium 2018, Dec. 8-9, 2018, Kawasaki}
\title{オープンCAEシンポジウム2018印刷用原稿の書き方(和文表題)}
\subtitle{\LaTeX のテンプレート(和文副題)}
\author{%
氏 名$^{1\dagger}$%
\hspace{1\zw}%
氏 名$^{2}$%
\hspace{1\zw}%
氏 名$^{3}$%
}
\affiliation{%
${}^{1}$所属%
\hspace{1\zw}%
${}^{2}$所属%
\hspace{1\zw}%
${}^{3}$所属%
} 
\etitle{Manuscripts preparation guide for the OpenCAE symposium 2018 (English Title)}
\esubtitle{The case of \LaTeX (English Sub-Title)}
\eauthor{%
Firstname FAMILYNAME$^{*}$%
\hspace{1em}%
Firstname FAMILYNAME$^{**}$%
\hspace{1em}%
Firstname FAMILYNAME$^{***}$%
}
\eaffiliation{%
${}^{*}$affiliation%
\hspace{1em}%
${}^{**}$affiliation%
\hspace{1em}%
${}^{***}$affiliation%
}
\keywords{Keyword1, Keyword2, Keyword3, Keyword4, Keyword5}
\AuthorsEmail{corresponding.author@opencae.or.jp}
\begin{abstract}
\end{abstract}
%
\begin{document}
\maketitle
%%%
\section{原稿について}
%%
\subsection{様式}
\begin{description}
\item [用紙]
  A4です．
\item [余白]
  上25mm，下25mm，右20mm，左20mmです．
\item [段組]
  1段組です．
\item [頁数]
  最小1ページ，最大{\bfseries 8ページ}です．
\end{description}
%%
\subsection{本文}
\begin{description}
\item [言語]
  日本語または英語です．
\item [句読点]
  本文が日本語の場合，
  句読点として，
  全角の読点「，」（カンマ）と句点「．」（ピリオド）を用いてください． 
\item [題目・所属]
  題目と所属については，以下とします．
  \begin{itemize}
  \item 
  本文が日本語の場合，日本語での題目，著者名と所属の記載に続けて，
  英語でも同内容を記載してください．
  \item 
  本文が英語の場合，英語の題目および著者名と所属のみの記載でも結構です．
  \end{itemize}
\item [キーワード]
  講演内容を良く表すキーワードを最低3語，通常5語程度選定し，英語で記入ください.
\item [字体]
  字体は，それぞれ以下とします．
  \begin{description}
  \item 
    [本文]
    明朝体・Serif系（Times New Romanなど）を使用してください．
  \item 
    [題目・所属・キーワード・見出し]
    ゴシック体・Sans-serif系（Helveticaなど）を使用してください．
  \end{description}
\item [文字の大きさ]
  フォントの大きさは，それぞれ以下とします．
  \begin{description}
  \item [題目]
    14ポイントです．
  \item [副題・大見出し]
    12ポイントです．
  \item [本文・所属・キーワード・図表キャプション・大見出し以外の見出し]
    10ポイントです．
  \end{description}
\item [参考文献]
  参考文献については，以下とします．
  \begin{itemize}
  \item 
  参考文献は，本文中の引用箇所の末尾に角括弧をつけた番号で表し，
  本文の末尾にまとめて列記してください．
  \item 
  WEBページについても，
  参考文献\cite{SIST02-200}等を参考にして，
  URLやアクセス日付を明記してください．
  \end{itemize}
\end{description}
%%
\subsection{図表}
\begin{description}
\item [言語]
  図表中の記号およびキャプションは英語とします．
\item [字体・文字の大きさ]
  明瞭である限り，図表中の字体や文字の大きさは任意です．
\item [本文中での参照]
  \autoref{fig:Logo}，
  \autoref{tab:Secretariat}などと記載してください．
\end{description}
%
\begin{figure}[htbp]
\centering
\subfigure[Color logo]{
\includegraphics[width=0.2\textwidth]{fig/Logo-c3-1.jpg}
\label{fig:Logo-c3-1}
}
\hspace{0.1\textwidth}
\subfigure[Monochrome logo]{
\includegraphics[width=0.2\textwidth]{fig/Logo-b3.jpg}
\label{fig:Logo-b3}
}
\caption{Logo of the OpenCAE Society}
\label{fig:Logo}
\end{figure}
%%%
\section{原稿提出について}
%%
\begin{description}
\item [ファイル形式]
  PDF形式で提出ください．
\item [ファイルサイズ]
  最大{\bfseries 20MB}です．
\item [フォント]
  提出されるPDFファイルに全てのフォントが埋め込れている事を確認ください．
\item [URL]
  できるだけURLにはリンクを付加してください．
\item [提出方法]
  \href
  {http://www.opencae.or.jp/activity/symposium/opencae_symposium2018/}
  {オープンCAEシンポジウム2018のWEBページ}
  をご参照の上，提出してください．
\end{description}
%%%
\section{お問い合わせ}
ご不明点については，
お手数ですが，
\autoref{tab:Secretariat}のシンポジウム事務局まで，
電子メールでお問い合わせください．
\begin{table}[htbp]
  \centering
  \begin{tabular}{l|l}
    \hline
    E-mail address & symposium2018@opencae.or.jp\\
    \hline
    \end{tabular}
  \caption{Secretariat of the OpenCAE symposium 2018}
  \label{tab:Secretariat}
\end{table}
%
\begin{thebibliography}{9}
\bibitem{SIST02-200}
科学技術情報流通技術基準
参照文献の書き方
SIST 02 – 2007,
独立行政法人科学技術振興機構(2007).
\href
{http://jipsti.jst.go.jp/sist/pdf/SIST02-2007.pdf}
{http://jipsti.jst.go.jp/sist/pdf/SIST02-2007.pdf}
(accessed 2015-09-29).
\end{thebibliography}
%
\end{document}

%% BiBTeXを用いる場合
%% template_OpenCAE_symposium.bibに参考文献データベースを記述し，
%% 以下の2行を有効にし，MakefileでBiBTeXを使うように設定してください．
%\bibliography{template_OpenCAE_symposium}
%\bibliographystyle{junsrt}
%%
%% 付録
\appendix
%%%
\section{コードの引用表示例}
%
コードの引用表示例を
\autoref{code:template_OpenCAE_symposium.tex}に示す．
%
\lstinputlisting[caption={Header of template\_OpenCAE\_symposium.tex.},language={tex},linerange={1-10},label={code:template_OpenCAE_symposium.tex}]
{./template_OpenCAE_symposium.tex}
%
\end{document}
